\catcode`\{=1 \catcode`\}=2 \catcode`\#=6
\font\tenrm=cmr10 \tenrm
\hsize=250pt \vsize=600pt \parindent=20pt \baselineskip=12pt \topskip=10pt
\tolerance=1000 \linepenalty=10 \maxdepth=2pt \parfillskip=0pt plus 1fil
\parskip=4pt plus 2pt
A first paragraph with the default parameters, long enough for two
lines or perhaps three at this measure.

\parindent=40pt \baselineskip=14pt \parskip=0pt
A second paragraph, indented more, with a larger baselineskip between
its lines than the first one had.

\noindent\hangindent=30pt \hangafter=1 A hanging paragraph, the lines after
the first indented by thirty points on the left as hangindent asks.

\hangindent=-30pt \hangafter=-2 A paragraph whose first two lines are
indented on the right, and the rest at the full measure, and on and on.

\parshape 3 10pt 200pt 20pt 180pt 0pt 250pt A parshape paragraph with
three line specifications given, the last one repeating for the rest of
the paragraph which goes on.

\leftskip=10pt \rightskip=10pt plus 20pt
A ragged paragraph with leftskip and rightskip set, which makes it ragged
right because the rightskip can stretch quite a lot.

\leftskip=0pt \rightskip=0pt
\looseness=1 A loose paragraph with looseness one, which TeX makes a line
longer if it can do so without exceeding the tolerance.

\emergencystretch=10pt \tolerance=100 \pretolerance=-1
An emergency paragraph with a low tolerance, where emergency stretch helps
the line breaker find a solution.

\tolerance=1000 \emergencystretch=0pt \pretolerance=0
\clubpenalty=150 \widowpenalty=150 \interlinepenalty=10 \brokenpenalty=100
Penalties between lines: club, widow and interline penalties appear
in the vertical list between the lines of this paragraph which is long.
\end
